\chapter{Computational Experiments} 

\label{part1:computationalexp}
In this section, we describe the results of our computational experiments with
our edge-based and set partitioning based formulations, and the valid inequalities.
%We ran our experiments on a 64-node Linux-based Beowulf cluster using condor. 

\section{Description of the Instances}
\label{part1:instances}

To create an instance, we need to specify customer and candidate facility locations, 
customer demands, vehicle capacity, vehicle time limit, facility capacity, 
and facility and vehicle fixed and variable costs. 

We used four MDVRP benchmark problems (called $p01$, $p03$, $p07$, $p11$) 
developed by \citet{Cordeau2} to generate customer locations and customer demands. 
Each of the \citet{Cordeau2} instances that we selected has a different set of customer 
locations. We did not select instance $p02$ since the instance has the same customer 
locations with the instance $p01$. Similarly, the instances $p04$, $p05$, $p06$
and $p07$ and the instances $p08$, $p09$, $p10$ and $p11$ have the same customer locations.


We generated 6 set of instances with 20, 25 and 30 customers. All of the instances have 
5 candidate facility locations. We included facility locations given for the related 
benchmark problem into the facility set and if the number of facilities is less than 
5 in the corresponding benchmark problem, we generated necessary number of facilities 
using a uniform distribution on the region where customers are located. 
We selected the facility capacity considering total demand and requiring at least 
two facilities to be in the solution. For each set of locations and demands, by using 2 
different vehicle capacity values and 2 different vehicle time limits, 
we generated 24 (4x6) instances with 20 customers, 24 instances with 25 customers
and 24 instances  with 30 customers. 

The smallest vehicle capacity for an instance was chosen a value that covers at least 
one customer or two customers and a customer that has the largest demand. The smallest 
vehicle time limit was calculated as a value that is greater than two times the distance 
between every customer and the furthest facility.

Table \ref{prt1:probparam} presents parameter values used and explains the generation of the instances.   
% the vehicle capacities, time limits and facility capacities 
%used for each group of instance.
\begin {table}[h]
\begin{center}
\caption {Generation of instances from \citet{Cordeau2} problems} \label{part1:probparam}
\begin {tabular}[t]{llll} 
Inst. Gorup	&	Benchmark Prb.	&	LRSP Instance 	&	Customer Locations and Demands	\\\hline
1	&	p01	&	p01-f	&	first 20, 25 and 30 customers	\\
2,3	&	p03	&	p03-f, p03-l	&	first and last 20, 25 and 30 customers	\\
4	&	p07	&	p07-f	&	first 20, 25 and 30 customers	\\
5,6	&	p11	&	p11-f, p11-l	&	first and last 20, 25 and 30 customers	\\\hline
\multicolumn{4}{c}{{\bf a. Customer information }} 
\end{tabular}
\begin {tabular}[t]{lccccc} 
%\multicolumn {6}{c}{{\bf Table 1.} Problem Parameters}\\ \hline
Inst. Generated Using & Vcap1	& Vcap2	& TimeLimit1 &	TimeLimit2 & Fcap\\\hline
p01,p07	& 	50	&	70	&	140	&	160	&	250 \\
p03	&	60	&	80	&	140	&	160	&	250\\
p11	&	150	&	200	&	240	&	260	&	500\\\hline
\multicolumn{6}{c}{{\bf b. Parameter information }} 
\end{tabular}
\end{center}
\end {table} 

We named our instances based on the parameters. For example, the instance 
$p01-f25-v1t1$ was created using the {\em first} 25 customer locations of  
the \citet{Cordeau2} instance called $p01$. The first possible vehicle capacity ($v1=50$) and time limit
($t1=140$) are used in the instance. Similarly, $p03-l25-v2t2$ was created using the {\em last} 
25 customer locations of the \citet{Cordeau2} instance called $p03$. The second possible vehicle 
capacity ($v2=80$) and time limit ($t2=160$) are used in the instance.  

To estimate the facility fixed cost, we used the study done by \citet{Burlett}. Assuming
a 100,000 sq. ft. facility, we estimated that the daily fixed cost of a facility varies 
between \$1500 and \$2300 depending on the region in which the facility is located. We 
randomly generated 5 values between 1500 and 2300. We used the same facility fixed 
costs in all of our instances. We estimated that the 
daily fixed cost of a mid-size vehicle varies between \$200 and \$250 based on the
economic life analysis done by  Kenworth Truck Company \citeyearpar{Kenworth}. We used
the average value \$225 as the fixed cost of a vehicle in all of our instances.  
The operating cost of a vehicle is estimated to be \$60 per hour. To estimate the
operating cost, we considered the average vehicle speed, the average amount of 
fuel consumption, fuel cost and the ratio of our working hour limit to the real
life working hour limit, which is generally 8 hours. Table \ref{costparam} lists 
the cost parameters used.  

\begin {table}[h]
\begin{center}
\caption {Cost Parameters} \label{costparam}
\begin {tabular}[t]{ccccccc} \hline
%\multicolumn {6}{c}{{\bf Table 1.} Problem Parameters}\\ \hline
$C^F_1$ & $C^F_2$ & $C^F_3$ & $C^F_4$ & $C^F_5$ & VF & OC \\\hline
\$1615 /day  &	\$1845 /day  &	\$1781 /day  &	\$1975 /day  &	\$1753 /day  &   \$225 /day  & \$ 60 /hour  	 \\\hline
\end{tabular}
\end{center}
\end {table} 

The travel time, $t_{ik}$, between each pair of nodes is calculated using 
Euclidean distance calculations and then rounded to the largest integer that is less
than or equal to the calculated value.  
We run the Floyd-Warshall Algorithm \citep{Ahuja} on the travel time
matrix to ensure that all travel times satisfy the triangle inequality.
In our test problems, we assume that the service time at each
node is zero. 

\section{Comparing the LP Relaxations of SPP-LRSP and EB-LRSP}
In section \ref{part1:dantzifW}, it is shown that the Dantzig Wolfe 
decomposition of edge-based model EB-LRSP is equal to the set 
partitioning based model, SPP-LRSP plus the bound 
constraint (\ref{conv2}) for the total number of vehicles 
located at each facility. In this section, we measure the strength 
of the LP relaxations of the model SPP-LRSP plus the upper bound constraint 
(\ref{conv2}) and the model EB-LRSP.   
We used instances with 25 customers and 30 customers in this experiment.
The LP relaxation of EB-LRSP was solved using CPLEX 9.1 and AMPL 
interface. The LP relaxation of SP-LRSP was solved using 
a column generation algorithm, which was coded in MINTO
(a mixed integer programming solver). We used CPLEX 9.1
as an LP solver in MINTO. All instances in both test groups 
were run at most 8 CPU hours.   

Based on Table (\ref{comparesp-eb-gr1}) and Table (\ref{comparesp-eb-gr2}), 
for all solved instances, the LP relaxations obtained from model SPP-LRSP are \%4.89 to 
\%12.82 higher than the LP relaxations obtained from model EB-LRSP.   
In addition, the computation times for the SPP-LRSP are much better than the times 
for EB-LRSP in all solved instances. Furthermore, only 2 instances cannot be
solved within time limit using model SPP-LRSP, while the number of unsolved instances
is 8 if model EB-LRSP is used. 

\begin{table}[!hbp]
\centering
\caption{Comparison of LP Relaxation Bounds of EB-LRS and SPP-LRS: 25 Customer Instances}\label{comparesp-eb-gr1}
\begin {tabular}[t]{|c|c|c|c|c|c|}\hline
\mbox{}	&  \multicolumn{2}{|c|}{{\bf SPP-LRS}} & \multicolumn{2}{|c|}{{\bf EB-LRS}}	&	{\bf (a-b)/b*100}	\\\cline{2-5}
{\bf Data File}	& {\bf LP Relax. (a)}	& {\bf CPU (s)} & {\bf LP Relax. (b)}	&  {\bf CPU (s)} & {\bf \% Gap}	\\\hline 
p01-f25-v1t1	&	3936.58	&	114.66	&	3705.06	&	5058.9	&	6.25	\\
p01-f25-v1t2	&	3855.91	&	346.18	&	3641.57	&	4719.67	&	5.89	\\
p01-f25-v2t1	&	3793.43	&	672.07	&	3610.19	&	4896.98	&	5.08	\\
p01-f25-v2t2	&	3724.91	&	4240.58	&	3551.19	&	5046.43	&	4.89	\\
p03-f25-v1t1	&	4620.03	&	26.18	&	4267.66	&	8969.62	&	8.26	\\
p03-f25-v1t2	&	4538.07	&	256.77	&	4195.35	&	8729.64	&	8.17	\\
p03-f25-v2t1	&	4443.34	&	121.03	&	4179.83	&	10256.6	&	6.3	\\
p03-f25-v2t2	&	4351.62	&	791.77	&	4113.72	&	10934.8	&	5.78	\\
p03-l25-v1t1	&	3928.73	&	34.83	&	3599.43	&	10906.3	&	9.15	\\
p03-l25-v1t2	&	3829.09	&	112.31	&	3527.39	&	10216.5	&	8.55	\\
p03-l25-v2t1	&	3812	&	78.6	&	3514.9	&	8722.87	&	8.45	\\
p03-l25-v2t2	&	3713.14	&	355.46	&	3445.92	&	8835.58	&	7.75	\\
p07-f25-v1t1	&	3413.54	&	63.77	&	3151.23	&	15809.4	&	8.32	\\
p07-f25-v1t2	&	3341.53	&	394.39	&	3085.95	&	12698.5	&	8.28	\\
p07-f25-v2t1	&	3301.11	&	198.84	&	3069.45	&	13462.4	&	7.55	\\
p07-f25-v2t2	&	3198.43	&	677.1	&	3007.23	&	11832.1	&	6.36	\\
p11-f25-v1t1	&	6615.99	&	0.42	&	5934.5	&	13097.1	&	11.48	\\
p11-f25-v1t2	&	6536.44	&	0.72	&	5865.22	&	14250.6	&	11.44	\\
p11-f25-v2t1	&	6378.77	&	0.73	&	5712.4	&	14334.6	&	11.67	\\
p11-f25-v2t2	&	6261.63	&	1.19	&	5650.04	&	14252.1	&	10.82	\\
p11-l25-v1t1	&	7298.5	&	0.26	&	6469.19	&	9429.75	&	12.82	\\
p11-l25-v1t2	&	7187.6	&	0.33	&	6402.79	&	9334.2	&	12.26	\\
p11-l25-v2t1	&	6939.53	&	0.36	&	6188.64	&	11521.1	&	12.13	\\
p11-l25-v2t2	&	6904.43	&	0.44	&	6131.83	&	11561.7	&	12.6	\\\hline
\end{tabular}
\end{table}

\begin{table}[!hbp]
\centering
\caption{Comparison of LP Relaxation Bounds of EB-LRS and SPP-LRS: 30 Customer Instances}\label{comparesp-eb-gr2}
\begin {tabular}[t]{|c|c|c|c|c|c|}\hline
\mbox{}	&  \multicolumn{2}{|c|}{{\bf SPP-LRS}} & \multicolumn{2}{|c|}{{\bf EB-LRS}}	&	{\bf (a-b)/b*100}	\\\cline{2-5}
{\bf Data File}	& {\bf LP Relax. (a)}	& {\bf CPU (s)} & {\bf LP Relax. (b)}	&  {\bf CPU (s)} & {\bf \% Gap}	\\\hline 
p01-f30-v1t1	&	4640.47	&	643.4	&	\mbox{}	&	\mbox{}	&	\mbox{}	\\
p01-f30-v1t2	&	4561.62	&	8721.2	&	4219.16	&	17751.3	&	8.12	\\
p01-f30-v2t1	&	4462.66	&	4464.09	&	4169.85	&	21347.8	&	7.02	\\
p01-f30-v2t2	&	\mbox{*}	&	29209.74	&	4106.6	&	18790	&	\mbox{-}	\\
p03-f30-v1t1	&	5413.72	&	370.4	&	\mbox{*}	&	28829.6	&	\mbox{-}	\\
p03-f30-v1t2	&	5330.33	&	6625.36	&	\mbox{*}	&	28827.6	&	\mbox{-}	\\
p03-f30-v2t1	&	5209.41	&	2953.18	&	\mbox{*}	&	28835	&	\mbox{-}	\\
p03-f30-v2t2	&	\mbox{*}	&	29356.69	&	\mbox{*}	&	28834.2	&	\mbox{-}	\\
p03-l30-v1t1	&	4707.72	&	131.18	&	4317.71	&	20964.7	&	9.03	\\
p03-l30-v1t2	&	4615.22	&	2806.27	&	4240.26	&	19245.8	&	8.84	\\
p03-l30-v2t1	&	4566.1	&	521.17	&	4221.57	&	19880.2	&	8.16	\\
p03-l30-v2t2	&	4459.03	&	8129.62	&	4148.77	&	19549.6	&	7.48	\\
p07-f30-v1t1	&	4153.16	&	705.12	&	3858.62	&	20347.3	&	7.63	\\
p07-f30-v1t2	&	4077.68	&	2625.43	&	3782.79	&	20362.7	&	7.8	\\
p07-f30-v2t1	&	3986.37	&	1870.87	&	3746.43	&	22084.3	&	6.4	\\
p07-f30-v2t2	&	3899.99	&	15385.17	&	3674.65	&	22243.4	&	6.13	\\
p11-f30-v1t1	&	7232.02	&	0.59	&	6483.4	&	17474.6	&	11.55	\\
p11-f30-v1t2	&	7124.96	&	0.9	&	6410.9	&	18558.7	&	11.14	\\
p11-f30-v2t1	&	6965.07	&	1.22	&	6238.88	&	22240	&	11.64	\\
p11-f30-v2t2	&	6875.87	&	1.81	&	6172.69	&	23976.6	&	11.39	\\
p11-l30-v1t1	&	8265.62	&	0.63	&	\mbox{*}	&	28843	&	\mbox{-}	\\
p11-l30-v1t2	&	8161.91	&	1	&	\mbox{*}	&	28823.6	&	\mbox{-}	\\
p11-l30-v2t1	&	8006.08	&	1.09	&	\mbox{*}	&	28817.9	&	\mbox{-}	\\
p11-l30-v2t2	&	7798.55	&	2.53	&	\mbox{*}	&	28816	&	\mbox{-}	\\\hline
\multicolumn{6}{l}{* Time limit is exceeded before a lp solution is found.}
\end{tabular}
\end{table}


\section{Comparison of Alternative Set Partitioning Formulations}

To compare the strength of the various set partitioning formulations obtained by
using the valid inequalities described in section \ref{StrongForm}, we 
enumerated all possible pairings and solved the formulations using a standard 
branch-and-bound algorithm. Since the number of feasible pairings is large, 
we are not able to enumerate all columns of the instances with 25 customers and 
30 customers. Instead, for the experiments in this part, we used the instances 
with 20 customers and 5 facilities. Furthermore, larger vehicle capacity and
time limit results more columns, therefore we only used the instances 
with small vehicle capacity and time limit (group v1t1). 
 
To generate all columns, we used an algorithm
written in C. We constructed mps files for formulations including generated 
columns using MINTO and solved the mps files using CPLEX 9.1.  
Since the constraints (\ref{spp_con1}) and 
(\ref{spp_con2}) are enough to determine the LRS problem, we construct our 
{\em base formulation ($SPP_B$)} from these constraints. 

{\bf Experiment 1} To compare the strength of constraints (\ref{relation-Z-T}) and 
(\ref{relation}) introduced in section \ref{StrongForm} and to verify the 
value of their contribution to the lower bound of the LRS problem,
we construct two formulations: $SPP_1$ and $SPP_2$. $SPP_1$ is constructed
by adding the constraints (\ref{relation-Z-T}) to the $SPP_B$ and $SPP_2$ 
is constructed by adding the constraints (\ref{relation}) to the $SPP_B$.
Table \ref{part1:formulations-gr1} shows the constraints (\ref{relation-Z-T}) and
(\ref{relation}), and lists the constraints in each formulation. 

\begin{table}[h]
\centering
\caption{LRS models for Experiment 1}\label{part1:formulations-gr1}
\begin {tabular}[t]{|l|l|l|}  
\multicolumn{3}{l}{\mbox{}}\\
\multicolumn{3}{l}{Const.(\ref{relation-Z-T}): $Z_{p} \leq T_{j} \ \ \forall j \in J, \forall p \in P_{j}$}\\
\multicolumn{3}{l}{\mbox{}}\\
\multicolumn{3}{l}{Const.(\ref{relation}): $\sum_{p \in P_{j}}a_{ip}Z_{p} \leq T_{j} \ \ \forall j \in J, \forall i \in N$}\\  
\multicolumn{3}{l}{\mbox{}}\\
\hline
$SPP_B$ & $SPP_1$ & $SPP_2$ \\\hline
1. SP Const.(\ref{spp_con1}) & 1. SP Const.(\ref{spp_con1}) & 1. SP Const.(\ref{spp_con1}) \\
2. Fac.Cap.Const.(\ref{spp_con2}) & 2. Fac.Cap.Const.(\ref{spp_con2}) & 2. Fac.Cap.Const.(\ref{spp_con2}) \\
\mbox{}  & 3. Relation Const.(\ref{relation-Z-T}) &  3. Relation Const.(\ref{relation})\\\hline
\end{tabular}
\end{table}

Table \ref{part1:lpbounds-gr1} compares the LP relaxation bounds of $SPP_B$, $SPP_1$ and $SPP_2$. In table 
\ref{part1:lpbounds-gr1}, \%Improv columns for models $SPP_1$ and $SPP_2$ represent the percentage increase
in the LP relaxation bound with respect to the LP relaxation bound of the base model $SPP_B$. The results 
suggest that 
the constraints (\ref{relation}) are better than the constraints (\ref{relation-Z-T}), since
the LP solution of the model $SPP_2$ is higher than the LP solution of the model $SPP_1$ in all
instances. The model $SPP_2$ improves the LP solution of the base model by 1.85\% to 8.34\% 
while the model $SPP_1$ only improves the LP solution of the base model by 0.01\% to 1.61\%.   

\begin{table}[h]
\centering
\caption{Results of Experiment 1}\label{part1:lpbounds-gr1}
\begin {tabular}[t]{|c|c|c|c|c|c|c|c|}  
\hline
{\bf Data File} & {\bf \# of Col.} & {\bf IP} & {\bf $SPP_B$}& \multicolumn{2}{|c|}{{\bf $SPP_1$}} & \multicolumn{2}{|c|}{{\bf $SPP_2$}}\\ \cline {4-8}
\mbox{} &  \mbox{} &  \mbox{} &  {\bf LP} & {\bf LP} & {\bf \%Improv.} & {\bf LP} & {\bf \%Improv.} \\\hline
p01-f20-v1t1	&	67553	&	4409	&	3335.23	&	3337.72	&	0.07	&	3545.31	&	6.3	\\
p03-f20-v1t1	&	52614	&	4459	&	3769.21	&	3769.55	&	0.01	&	3838.82	&	1.85	\\
p03-l20-v1t1	&	32796	&	4610	&	3173.4	&	3192.92	&	0.62	&	3416.8	&	7.67	\\
p07-f20-v1t1	&	68262	&	4420	&	2734.17	&	2745.73	&	0.42	&	2886.51	&	5.57	\\
p11-f20-v1t1	&	1945	&	6129	&	5594.08	&	5625.36	&	0.56	&	6060.6	&	8.34	\\
p11-f20-v1t1	&	861	&	7845	&	5988.01	&	6084.28	&	1.61	&	6420.39	&	7.22	\\\hline
\end{tabular}
\end{table} 

\begin{comment} %%%%%%%%% PREVIOUS COMPUTATIONAL RESULTS %%%%%%%%%%%%%%%%%%%%%%%%%%%%%%%%%%%%%%%
\begin{table}[h]
\centering
\caption{Results of Experiment 1}\label{lpbounds-gr1}
\begin {tabular}[t]{|c|c|c|c|c|c|c|c|}  
\hline
{\bf Data File} & {\bf \# of Col.} & {\bf IP} & {\bf $SPP_B$}& \multicolumn{2}{|c|}{{\bf $SPP_1$}} & \multicolumn{2}{|c|}{{\bf $SPP_2$}}\\ \cline {4-8}
\mbox{} &  \mbox{} &  \mbox{} &  {\bf LP} & {\bf LP} & {\bf \%Improv.} & {\bf LP} & {\bf \%Improv.} \\\hline
p01-f20	&	67558	&	4348	&	3273.822	&	3276.9	&	0.09	&	3471.73	&	6.045	\\\hline
p03-f20	&	52619	&	4394.33	&	3705.03	&	3705.07	&	0	&	3767.8	& 1.69		\\\hline
p03-l20	&	14779	&	4708.83	&	3552.272	&	3552.61	&	0.01	&	3755.31	&	5.716	\\\hline
p07-f20	&	68267	&	4357.17	&	2677.031	&	2687.63	&	0.4	&	2813.44	&	5.095	\\\hline
p11-f20	&	1950	&	5921.5	&	5291.045	&	5308.24	&	0.32	&	5849.5	&	10.555	\\\hline
p11-l20	&	898	&	8017.5	&	5865.493	&	5987.21	&	2.08	&	6855.28	&	16.875	\\\hline
\end{tabular}
\end{table} 
\end{comment}


{\bf Experiment 2} After we observing that constraints (\ref{relation}) seem to be better than constraints 
(\ref{relation-Z-T}), we design experiment 2 to test the strength of constraint (\ref{lb-open-fac}) 
with respect to the base model and to test the combined affect of constraints (\ref{lb-open-fac}) and (\ref{relation}).
We construct two formulations: $SPP_3$ and $SPP_4$. $SPP_3$ is constructed
by adding the constraint (\ref{lb-open-fac}) to the base model $SPP_B$. 
$SPP_4$ is constructed by adding the constraints (\ref{relation}) to the $SPP_3$.
Table \ref{part1:formulations-gr2} shows the constraint (\ref{lb-open-fac}) and lists the constraints in each formulation.   

\begin{table}[h]
\centering
\caption{LRS models for Experiment 2}\label{part1:formulations-gr2}
\begin {tabular}[t]{|l|l|l|}  
\multicolumn{3}{l}{\mbox{}}\\
\multicolumn{3}{l}{Const.(\ref{lb-open-fac}): $\sum_{j \in J} T_{j} \geq nREQ$ }\\
\multicolumn{3}{l}{\mbox{}}\\
\hline
$SPP_B$ & $SPP_3$ & $SPP_4$ \\\hline
1. SP Cont.(\ref{spp_con1}) & 1. SP Cont.(\ref{spp_con1}) & 1. SP Cont.(\ref{spp_con1}) \\
2. Fac.Cap.Const.(\ref{spp_con2}) & 2. Fac.Cap.Const.(\ref{spp_con2}) & 2. Fac.Cap.Const.(\ref{spp_con2}) \\
\mbox{} & 3. Const.(\ref{lb-open-fac})  & 3. Const.(\ref{lb-open-fac}) \\
\mbox{} & \mbox{}  & 4. Relation Cont.(\ref{relation})\\\hline
\end{tabular}
\end{table}

Table \ref{part1:lpbounds-gr2} presents the percentage improvement in the LP relaxations with
constraint (\ref{lb-open-fac}) and the percentage improvement in the LP relaxations with
constraints (\ref{lb-open-fac}) and (\ref{relation}). In table \ref{pat1:lpbounds-gr2}, the \%Improv. 
columns for models $SPP_2$, $SPP_3$ and $SPP_4$ represent the percentage increase
in the LP solutions with respect to the LP solution of the base model $SPP_B$. 
The results show that the LP bound improves by 5.39\% to 57.33\% with constraint 
(\ref{lb-open-fac}). When we compare models $SPP_2$ and $SPP_3$, we observe that
in most instances, constraint (\ref{lb-open-fac}) improves the LP relaxation bound 
more than the constraints (\ref{relation}) do. The two set of constraints, 
constraints (\ref{relation}) and
constraint (\ref{lb-open-fac}) together achieve larger improvement than 
constraint (\ref{lb-open-fac}) does alone in all instances. The LP relaxation
solution of the model $SPP_4$ is 8.34\% to 60.05\% better than the LP
solution of the base model. 

\begin{sidewaystable}
%\begin{table}[h]
\centering
\caption{Results of Experiment 2}\label{part1:lpbounds-gr2}
\begin {tabular}[t]{|c|c|c|c|c|c|c|c|c|c|}  
\hline
{\bf Data File} & {\bf \# of Col.} & {\bf IP} & {\bf $SPP_B$}& \multicolumn{2}{|c|}{{\bf $SPP_2$}} & \multicolumn{2}{|c|}{{\bf $SPP_3$}} & \multicolumn{2}{|c|}{{\bf $SPP_4$}}\\ \cline {4-10}
\mbox{} &  \mbox{} &  \mbox{} &  {\bf LP} & {\bf LP} & {\bf \%Improv.} & {\bf LP} & {\bf \%Improv.} & {\bf LP} & {\bf \%Improv.} \\\hline
p01-f20-v1t1	&	67553	&	4409	&	3335.23		& 3545.31	& 6.3 & 4310.45 & 29.24	& 4341.31	& 30.17	\\
p03-f20-v1t1	&	52614	&	4459	&	3769.21		&	3838.82	&	1.85	&	4436.01	&	17.69	&	4448.7	&	18.03	\\
p03-l20-v1t1	&	32796	&	4610	&	3173.4		&	3416.8	&	7.67	&	4398.3	&	38.6	&	4537.44	&	42.98	\\
p07-f20-v1t1	&	68262	&	4420	&	2734.17		&	2886.51	&	5.57	&	4301.54	&	57.33	&	4376.1	&	60.05	\\
p11-f20-v1t1	&	1945	&	6129	&	5594.08		&	6060.6	&	8.34	&	5895.79	&	5.39	&	6060.6	&	8.34	\\
p11-f20-v1t1	&	861	&	7845	&	5988.01		&	6420.39	&	7.22	&	7545.86	&	26.02	&	7698.17	&	28.56	\\\hline
\end{tabular}
%\end{table} 
\end{sidewaystable}


\begin{comment} %%%%%%%%% PREVIOUS COMPUTATIONAL RESULTS %%%%%%%%%%%%%%%%%%%%%%%%%%%%%%%%%%%%%%%
\begin{sidewaystable}
%\begin{table}[h]
\centering
\caption{Results of Experiment 2}\label{lpbounds-gr2}
\begin {tabular}[t]{|c|c|c|c|c|c|c|c|c|c|}  
\hline
{\bf Data File} & {\bf \# of Col.} & {\bf IP} & {\bf $SPP_B$}& \multicolumn{2}{|c|}{{\bf $SPP_2$}} & \multicolumn{2}{|c|}{{\bf $SPP_3$}} & \multicolumn{2}{|c|}{{\bf $SPP_4$}}\\ \cline {4-10}
\mbox{} &  \mbox{} &  \mbox{} &  {\bf LP} & {\bf LP} & {\bf \%Improv.} & {\bf LP} & {\bf \%Improv.} & {\bf LP} & {\bf \%Improv.} \\\hline
p01-f20	&	67558	&	4348	&	3273.822        & 3471.73	& 6.045  & 4251.74	&	29.87	&	4279.88	&	30.73	\\\hline
p03-f20	&	52619	&	4394.33	&	3705.03	        & 3767.8	& 1.69   &  4371.08	&	17.98	&	4383.24	&	18.31	\\\hline
p03-l20	&	14779	&	4708.83	&	3552.272	& 3755.31	& 5.716  & 4514.81	&	27.1	&	4656.04	&	31.07	\\\hline
p07-f20	&	68267	&	4357.17	&	2677.031	& 2813.44	& 5.095  & 4243.97	&	58.53	&	4312.98	&	61.11	\\\hline
p11-f20	&	1950	&	5921.5	&	5291.045	& 5849.5	& 10.555 & 5623.73	&	6.29	&	5849.5	&	10.55	\\\hline
p11-l20	&	898	&	8017.5	&	5865.493	& 6855.28	& 16.875 & 7198.17	&	22.72	&	7810.83	&	33.17	\\\hline
\end{tabular}
%\end{table} 
\end{sidewaystable}
\end{comment}

Based on the experiments 1 and 2, we conclude that our valid inequalities improve the LP relaxation bound of the model, 
constraints (\ref{relation}) are better than  constraints (\ref{relation-Z-T}) and the $SPP_4$ formulation, which includes 
constraints (\ref{lb-open-fac}) and (\ref{relation}), is the best among other models since it has the highest 
lower bound in all instances. 


\qquad
















\begin{comment}
In chapter \ref{computationalexp}, we empirically demonstrate the improvement in the LP relaxation 
bounds achieved by constraints (\ref{lb-open-fac}), (\ref{relation-Z-T}) and (\ref{relation}). 
Throughout the remainder of the proposal, we use the augmented SPP-LRS formulation (AUG-SPP-LRS),
which includes the objective function (\ref{spp-obj}) and the constraints (\ref{spp_con1}), 
(\ref{spp_con2}), (\ref{binary1}), (\ref{binary2}), (\ref{lb-open-fac}), and (\ref{relation}). 
\end{comment}

